

\chapter*{Introduction Générale}
\markboth{\MakeUppercase{Introduction Générale}}{}%
\addcontentsline{toc}{chapter}{Introduction Générale}%

\vspace{1cm}

L'avènement des modèles de langage de grande taille (Large Language Models - LLM) a révolutionné le domaine de l'intelligence artificielle et ouvert de nouvelles perspectives dans le secteur de l'éducation et de l'apprentissage. Des technologies comme GPT-4 d'OpenAI, Claude d'Anthropic et Gemini de Google démontrent des capacités remarquables de compréhension et de génération de texte, offrant ainsi des opportunités inédites pour l'accompagnement pédagogique et la recherche d'information.

Dans ce contexte d'innovation technologique, les apprenants et professionnels cherchent des outils capables d'exploiter ces avancées tout en offrant une expérience utilisateur fluide et intuitive. La multiplicité des modèles disponibles pose néanmoins un défi : comment tirer profit des forces spécifiques de chaque modèle sans avoir à naviguer entre plusieurs plateformes distinctes ?

\section*{Problématique}

Face à la diversité croissante des modèles d'IA et à la fragmentation des outils d'apprentissage en ligne, plusieurs problèmes se posent :

\begin{itemize}
    \item L'accès aux différents modèles de langage (GPT-4, Claude, Gemini) nécessite l'utilisation de multiples plateformes avec des abonnements séparés
    \item Le manque d'outils intégrés pour l'analyse de documents, la génération de contenu visuel et l'accès à des ressources complémentaires
    \item L'absence d'espaces collaboratifs permettant aux utilisateurs de partager leurs connaissances et expériences
    \item La difficulté à générer des exercices personnalisés adaptés aux besoins spécifiques des apprenants
    \item Le besoin d'une interface unifiée et intuitive facilitant l'interaction avec ces technologies complexes
\end{itemize}

\section*{Objectifs du Projet}

L'objectif principal de ce projet est de concevoir et développer \textbf{SynapseAI}, une plateforme intelligente qui répond aux défis énoncés en proposant :

\begin{enumerate}
    \item \textbf{Interface Multi-LLM Unifiée} : Un accès centralisé à plusieurs modèles d'intelligence artificielle (GPT-4, Claude, Gemini) via une interface unique
    \item \textbf{Analyse Intelligente de Documents} : Un système permettant d'extraire les concepts clés et de générer des résumés à partir de documents uploadés
    \item \textbf{Génération de Graphiques} : Un outil de création automatique de visualisations (graphiques linéaires, à barres, circulaires) basé sur Chart.js
    \item \textbf{Recommandations YouTube Contextuelles} : Un système intelligent qui suggère des vidéos éducatives pertinentes en fonction des sujets de conversation
    \item \textbf{Espace Communautaire} : Une plateforme sociale permettant la publication de posts, l'ajout de commentaires et l'interaction entre utilisateurs
    \item \textbf{Générateur d'Exercices} : Un système de création d'exercices personnalisés avec correction automatique
    \item \textbf{Gestion Administrative} : Un tableau de bord complet pour la supervision des utilisateurs et du contenu de la plateforme
\end{enumerate}

\section*{Méthodologie de Travail}

Pour mener à bien ce projet, nous avons adopté la méthodologie agile \textbf{SCRUM} qui permet une approche itérative et incrémentale du développement. Le projet a été décomposé en trois sprints principaux :

\begin{itemize}
    \item \textbf{Sprint 1} : Mise en place de l'architecture de base, authentification des utilisateurs, intégration des modèles LLM et interface de chat
    \item \textbf{Sprint 2} : Développement des fonctionnalités avancées (analyse de documents, génération de graphiques, recommandations YouTube)
    \item \textbf{Sprint 3} : Implémentation de l'espace communautaire, du générateur d'exercices et du tableau de bord administratif
\end{itemize}

\section*{Organisation du Rapport}

Ce rapport est organisé en cinq chapitres principaux :

\begin{itemize}
    \item Le \textbf{Chapitre 1} présente le contexte général du projet, la problématique, l'étude de l'existant et la méthodologie SCRUM adoptée
    \item Le \textbf{Chapitre 2} expose l'état de l'art des technologies utilisées, la spécification des besoins fonctionnels et non fonctionnels, ainsi que la modélisation UML
    \item Le \textbf{Chapitre 3} détaille la réalisation du Sprint 1 incluant l'authentification, l'intégration multi-LLM et l'interface de chat
    \item Le \textbf{Chapitre 4} décrit le Sprint 2 avec l'implémentation de l'analyse de documents, la génération de graphiques et les recommandations YouTube
    \item Le \textbf{Chapitre 5} présente le Sprint 3 avec l'espace communautaire, le générateur d'exercices et le tableau de bord administratif
\end{itemize}

Enfin, une conclusion générale synthétise les réalisations du projet et propose des perspectives d'évolution de la plateforme SynapseAI.

